\chapter{Chapter 1 -- Medical Context and AI Foundations for Healthcare Data Analysis}
\section{Introduction}


\section{Healthcare Data and Collaborating Institutions}

Healthcare data is one of the most valuable and sensitive types of information in modern society, as it contains detailed records about patients’ physiological conditions, diagnoses, treatments, and outcomes \cite{rumsfeld2016big}. These data are generated and collected across a wide range of healthcare institutions, including hospitals, specialized clinics, general practitioner offices, and research centers. Each institution typically maintains its own information systems and databases, resulting in fragmented data storage and limited interoperability across institutions \cite{jensen2012mining}.

\subsection{Healthcare Institutions and Data Silos}

Healthcare institutions differ in scale, specialization, and technological infrastructure. Large hospitals often rely on Electronic Health Records (EHRs), which store structured data such as laboratory results and prescriptions, as well as unstructured data such as clinical notes and medical reports \cite{jensen2012mining}. In contrast, specialized centers may generate high-dimensional data such as medical imaging or physiological signals like electroencephalography (EEG).

This heterogeneity contributes to the emergence of \textit{data silos}, which refer to isolated repositories of data that are not easily accessible or shareable across systems or institutions. Data silos arise due to differences in database formats, lack of interoperability standards, and institutional policies restricting data access. As a result, the potential of large-scale data analysis and machine learning is significantly limited, since robust models require diverse and representative datasets \cite{mandl2001public}.

\subsection{Challenges of Data Sharing in Healthcare}

Sharing healthcare data across institutions is a complex process that involves multiple challenges. First, there are \textbf{technical challenges}, including heterogeneous data formats, inconsistent data quality, and lack of standardized terminologies. For instance, similar clinical variables may be recorded differently across institutions, making integration difficult \cite{kahn2016transparent}.

Second, \textbf{organizational challenges} play a major role. Healthcare institutions may be reluctant to share data due to concerns about data ownership, competitive advantage, or administrative complexity. Establishing agreements for data sharing often requires significant coordination and trust between stakeholders.

Third, \textbf{privacy and security concerns} are critical. Healthcare data contains highly sensitive personal information, and any breach or misuse can have severe consequences for patients. Therefore, strict measures must be taken to ensure data confidentiality, integrity, and secure access \cite{shickel2017deep}.

\subsection{Legal and Ethical Constraints (GDPR, HIPAA)}

The sharing and processing of healthcare data are governed by strict legal and ethical regulations. In the European context, the \textit{General Data Protection Regulation (GDPR)} establishes comprehensive rules for the protection of personal data. It requires explicit patient consent, enforces data minimization principles, and mandates anonymization or pseudonymization when data is used for research purposes \cite{EU2016GDPR}.

Similarly, in the United States, the \textit{Health Insurance Portability and Accountability Act (HIPAA)} defines standards for safeguarding medical information. HIPAA regulates how healthcare providers and organizations store, access, and transmit patient data, and it imposes strict penalties in case of data breaches \cite{hipaa1996}.

These regulations highlight the importance of developing methods that enable data utilization while preserving patient privacy.

\subsection{Need for Collaborative Learning in Healthcare}

Given the challenges associated with data sharing, there is a growing need for approaches that allow collaborative analysis without requiring direct data exchange. Traditional centralized machine learning approaches, which rely on aggregating all data into a single repository, are often impractical in healthcare due to regulatory and privacy constraints.

As a result, distributed learning paradigms have emerged as a promising alternative. These approaches enable multiple institutions to collaboratively train models while keeping the data locally stored. Such strategies allow leveraging large-scale, diverse datasets while preserving data privacy, thereby addressing one of the key limitations of conventional machine learning in healthcare \cite{kairouz2021advances}.



\section{Medical Data Modalities}

Healthcare data is inherently heterogeneous, encompassing multiple modalities that differ in structure, dimensionality, and acquisition processes. These modalities include physiological signals, medical imaging, and structured clinical data, each providing complementary information about patient health. The integration of such diverse data sources is essential for accurate diagnosis, prognosis, and personalized treatment, but it also introduces significant computational and methodological challenges \cite{esteva2019guide}.

\subsection{EEG Signals}

\subsubsection{Definition and Characteristics}

Electroencephalography (EEG) is a non-invasive technique used to record electrical activity generated by the brain through electrodes placed on the scalp. EEG signals are time-series data characterized by high temporal resolution, typically measured in milliseconds, which makes them particularly suitable for capturing dynamic neural processes \cite{niedermeyer2005eeg}.

EEG signals are commonly analyzed in both the time and frequency domains, as brain activity exhibits oscillatory patterns that can be decomposed into distinct frequency bands. These oscillations reflect different physiological and cognitive states, and their analysis provides valuable insights into brain function \cite{lopes2010eeg, eeg2015bands}.

EEG signals are typically categorized into several frequency bands, each associated with specific neurological activities:

\begin{itemize}
    \item \textbf{Delta waves} (0.5–4 Hz): primarily associated with deep sleep and unconscious states \cite{niedermeyer2005eeg}.
    \item \textbf{Theta waves} (4–8 Hz): linked to drowsiness, meditation, and memory-related processes \cite{klimesch1999eeg}.
    \item \textbf{Alpha waves} (8–13 Hz): dominant during relaxed wakefulness, particularly with closed eyes \cite{eeg2015bands}.
    \item \textbf{Beta waves} (13–30 Hz): associated with active thinking, concentration, and problem-solving \cite{pfurtscheller1999event}.
    \item \textbf{Gamma waves} (>30 Hz): related to higher cognitive functions such as perception and consciousness \cite{buzsaki2006rhythms}.
\end{itemize}

The analysis of these frequency bands is particularly important in clinical contexts, as abnormal patterns in specific bands can indicate neurological disorders. In the case of epilepsy, for example, characteristic disruptions and spikes in EEG signals can be observed across different frequency ranges, making spectral analysis a key tool for detection and diagnosis.

However, EEG signals are often noisy, non-stationary, and highly sensitive to external artifacts such as muscle movements and environmental interference, making their analysis challenging.

\subsubsection{EEG in Neurological Disease Monitoring}

EEG plays a critical role in the monitoring and diagnosis of neurological disorders. It is widely used to study brain activity patterns associated with conditions such as epilepsy, sleep disorders, and brain injuries. Machine learning and deep learning techniques have increasingly been applied to EEG data to automate the detection of abnormal patterns and improve diagnostic accuracy \cite{roy2019deep}. These approaches leverage the temporal structure of EEG signals to identify subtle variations that may not be easily detectable by clinicians.

\subsubsection{Epilepsy: Definition, Types, and Clinical Context}

Epilepsy is a chronic neurological disorder characterized by recurrent, unprovoked seizures resulting from abnormal electrical activity in the brain. It affects millions of individuals worldwide and presents in various forms, including focal seizures and generalized seizures \cite{who2019epilepsy}. EEG remains a primary tool for epilepsy diagnosis, as it allows clinicians to observe characteristic patterns such as spikes and wave discharges. The complexity and variability of EEG signals in epileptic patients make this domain particularly suitable for advanced machine learning applications.

\subsection{Medical Imaging}

\subsubsection{Overview of Medical Imaging Modalities}

Medical imaging provides visual representations of internal body structures and is a cornerstone of modern diagnostic medicine. Common imaging modalities include X-ray, Computed Tomography (CT), Magnetic Resonance Imaging (MRI), and ultrasound, each offering distinct advantages in terms of resolution, contrast, and clinical applicability \cite{litjens2017survey}. These imaging techniques generate high-dimensional data that require specialized processing methods for analysis.

\subsubsection{Cancer Sub-Types and Imaging-Based Diagnosis}

Medical imaging is extensively used in oncology for tumor detection, classification, and monitoring. Different cancer sub-types exhibit distinct visual patterns in imaging data, which can be exploited using machine learning and deep learning models for automated diagnosis. Convolutional Neural Networks (CNNs), in particular, have demonstrated remarkable performance in image-based cancer classification tasks, enabling early detection and improved clinical decision-making \cite{esteva2017dermatologist}. However, variations in imaging protocols and equipment across institutions can introduce significant variability in the data.

\subsection{Clinical and Structured Patient Data}

In addition to signals and images, healthcare systems generate structured data in the form of electronic health records. These include demographic information, laboratory results, medication history, and clinical observations. Such data are typically stored in tabular formats and are essential for predictive modeling tasks such as disease risk assessment and outcome prediction \cite{miotto2017deep}. Despite being more structured, clinical data often suffer from missing values, inconsistencies, and biases, which must be addressed during preprocessing.

\subsection{Challenges of Multimodal Medical Data Integration}

The integration of multiple medical data modalities presents significant challenges. First, modalities differ in terms of scale and representation, ranging from time-series signals (EEG) to high-dimensional images and structured tabular data. Second, aligning these modalities requires synchronization and proper feature extraction techniques. Third, multimodal datasets are often incomplete, as not all patients undergo the same tests or imaging procedures.

From a machine learning perspective, combining multimodal data can improve model performance by capturing complementary information. However, it also increases model complexity and computational cost. These challenges are further amplified in distributed settings, where data are stored across multiple institutions, reinforcing the need for advanced learning paradigms capable of handling multimodal and decentralized data \cite{ngiam2011multimodal}.

\section{Artificial Intelligence in Healthcare}

Artificial intelligence (AI) refers to the branch of computer science that aims to develop systems capable of performing tasks that typically require human intelligence, such as reasoning, perception, and decision-making \cite{russell2016ai}. In healthcare, AI has shown significant potential to improve disease diagnosis, prognosis, and patient management by leveraging large and complex medical datasets.

\subsection{Machine Learning}

\subsubsection{Definition}

Machine learning (ML) is a subfield of AI in which algorithms automatically learn patterns from data to make predictions or decisions without explicit programming \cite{mitchell1997machine}. ML models are trained using historical patient data, such as clinical records, laboratory tests, or medical imaging, and can then generalize to unseen cases. In healthcare, ML has been applied for tasks such as disease classification, risk prediction, and patient outcome forecasting.

\subsubsection{Types of Machine Learning}

\paragraph{Supervised Learning}  
Supervised learning involves training a model on labeled data, where each input is associated with a known output or target variable. The model learns a mapping function $f: X \rightarrow Y$, where $X$ represents input features and $Y$ the target labels. Examples in healthcare include tumor classification from imaging data and predicting patient readmission rates \cite{esteva2017dermatologist}.

\paragraph{Unsupervised Learning}  
Unsupervised learning deals with unlabeled data, where the algorithm identifies patterns or groupings within the dataset. Common methods include clustering and dimensionality reduction. In healthcare, unsupervised learning can uncover hidden subgroups of patients, aiding in personalized treatment strategies \cite{hinton2006reducing}.

\paragraph{Semi-Supervised Learning}  
Semi-supervised learning is used when labeled data is limited but unlabeled data is abundant. This approach is particularly relevant in medical applications, where annotating data (e.g., segmenting tumors) is costly and time-consuming \cite{chapelle2009semi}.

\paragraph{Reinforcement Learning}  
Reinforcement learning (RL) involves an agent that interacts with an environment to learn optimal policies through trial and error. In healthcare, RL has been explored for sequential treatment planning and personalized therapy optimization \cite{gottesman2019guidelines}.

\subsubsection{Popular Machine Learning Algorithms}

\begin{itemize}
    \item \textbf{Logistic Regression}: Used for binary or multi-class classification; predicts the probability of a class using a logistic function.
    \item \textbf{Decision Trees}: Hierarchical structures that split data based on feature values; interpretable and widely used in clinical decision support.
    \item \textbf{Random Forest}: Ensemble of decision trees; improves generalization and reduces overfitting.
    \item \textbf{Support Vector Machines (SVM)}: Finds the hyperplane that maximally separates classes in feature space.
    \item \textbf{k-Nearest Neighbors (k-NN)}: Classifies samples based on the majority label of their k nearest neighbors.
    \item \textbf{Gradient Boosting / XGBoost}: Ensemble methods that sequentially train weak learners to correct previous errors.
\end{itemize}

\subsubsection{Challenges of Machine Learning in Healthcare}

Healthcare data presents several challenges for ML, including heterogeneity of data modalities, class imbalance (e.g., rare diseases), missing values, and the need for interpretability and privacy preservation \cite{miotto2018deep}.

\subsection{Deep Learning}

\subsubsection{Definition and Principles}

Deep learning (DL) is a subset of machine learning that uses artificial neural networks with multiple layers to automatically learn hierarchical feature representations from raw data \cite{lecun2015deep}. A typical neural network consists of an input layer, one or more hidden layers, and an output layer. Each neuron applies a weighted sum of its inputs followed by a non-linear activation function $f$:

\[
y = f\left(\sum_{i=1}^{n} w_i x_i + b\right)
\]

where $w_i$ are weights, $x_i$ are inputs, $b$ is the bias, and $f$ is an activation function such as ReLU, sigmoid, or softmax.

\subsubsection{Artificial Neural Networks (ANNs)}

ANNs are the basic building blocks of deep learning. They can approximate complex, non-linear functions and are trained using backpropagation, which computes gradients of the loss function with respect to weights and updates them using an optimizer such as stochastic gradient descent (SGD) or Adam \cite{kingma2014adam}.

\subsubsection{Convolutional Neural Networks (CNNs)}

CNNs are specialized neural networks for grid-like data, such as images. Key components include:

\begin{itemize}
    \item \textbf{Convolutional layers}: Apply filters (kernels) to extract spatial features.
    \item \textbf{Pooling layers}: Reduce spatial dimensions, retaining essential information.
    \item \textbf{Fully connected layers}: Combine extracted features for final classification.
\end{itemize}

CNNs have been successfully applied in radiology for detecting tumors, lesions, and other abnormalities in imaging data \cite{litjens2017survey}.

\subsubsection{Recurrent Neural Networks (RNNs)}

RNNs are designed for sequential data, where each output depends on previous computations. Long Short-Term Memory (LSTM) and Gated Recurrent Unit (GRU) networks address long-term dependencies, making them suitable for EEG and ECG signal analysis \cite{lipton2015critical}.

\subsubsection{Deep Learning for Medical Data}

Deep learning can process diverse medical modalities, including imaging (CNNs), temporal signals like EEG/ECG (RNNs), and multimodal fusion combining clinical, imaging, and genomic data \cite{miotto2018deep}. DL models can automatically learn hierarchical features, reducing the need for manual feature engineering.

\subsection{Model Evaluation Metrics}

\subsubsection{Classification Metrics}

\begin{itemize}
    \item \textbf{Accuracy}: Fraction of correctly predicted instances.  
    \item \textbf{Precision}: Fraction of true positives among predicted positives.  
    \item \textbf{Recall (Sensitivity)}: Fraction of true positives among actual positives.  
    \item \textbf{F1-Score}: Harmonic mean of precision and recall:
    \[
    F1 = 2 \times \frac{\text{Precision} \times \text{Recall}}{\text{Precision + Recall}}
    \]
\end{itemize}

\subsubsection{Confusion Matrix}

The confusion matrix provides a complete summary of prediction performance, showing true positives (TP), false positives (FP), true negatives (TN), and false negatives (FN).

\subsubsection{ROC Curve and AUC}

The Receiver Operating Characteristic (ROC) curve plots the true positive rate against the false positive rate at various thresholds. The area under the curve (AUC) measures the classifier’s discrimination ability.

\subsubsection{Optimizers and Loss Functions}

Common optimizers include stochastic gradient descent (SGD), Adam, and RMSProp. Loss functions quantify the model error; for classification, cross-entropy loss is widely used:

\[
L = - \sum_{i=1}^{N} y_i \log(\hat{y}_i)
\]

where $y_i$ is the true label and $\hat{y}_i$ is the predicted probability.

\subsubsection{Overfitting and Regularization}

Overfitting occurs when a model memorizes training data but fails to generalize. Regularization techniques include L1/L2 penalties, dropout, and early stopping.

% ============================================================
\section{Artificial Intelligence in Healthcare}
% ============================================================

Artificial intelligence (AI) refers to the branch of computer science that aims to develop 
systems capable of performing tasks that typically require human intelligence, such as 
reasoning, perception, and decision-making \cite{russell2016ai}. In healthcare, AI has shown 
significant potential to transform disease diagnosis, prognosis, and patient management by 
leveraging large and complex medical datasets — including clinical records, medical images, 
and biosignals such as electroencephalograms (EEG). The following sections present the 
foundational AI concepts that underpin this work, from classical machine learning to deep 
learning architectures, along with the evaluation framework used to assess model performance.

% ============================================================
\subsection{Machine Learning}
% ============================================================

\subsubsection{Definition}

Machine learning (ML) is a subfield of AI in which algorithms automatically learn patterns 
from data to make predictions or decisions without being explicitly programmed for each 
task \cite{mitchell1997machine}. ML models are trained on historical data — such as patient 
clinical records, laboratory results, or annotated medical images — and are then expected 
to generalize to unseen cases. In healthcare, ML has been applied to a wide range of 
problems including disease classification, survival prediction, patient stratification, 
and clinical decision support.

% ------------------------------------------------------------
\subsubsection{Types of Machine Learning}
% ------------------------------------------------------------

\paragraph{Supervised Learning}
Supervised learning involves training a model on a labeled dataset, where each input sample 
$\mathbf{x} \in \mathcal{X}$ is paired with a known output $y \in \mathcal{Y}$. The model 
learns a mapping function $f: \mathcal{X} \rightarrow \mathcal{Y}$ by minimizing a loss 
function over the training data. In healthcare, supervised learning is the most widely 
used paradigm, with applications including tumor classification from imaging data, 
epilepsy detection from EEG signals, and prediction of patient readmission 
rates \cite{esteva2017dermatologist}.

\paragraph{Unsupervised Learning}
Unsupervised learning deals with unlabeled data, where the algorithm seeks to discover 
hidden structure or groupings within the dataset. Common techniques include clustering 
(e.g., K-means, DBSCAN) and dimensionality reduction (e.g., PCA, autoencoders). In 
healthcare, unsupervised methods are used to identify patient subgroups, uncover latent 
disease phenotypes, and support exploratory data analysis \cite{hinton2006reducing}.

\paragraph{Semi-Supervised Learning}
Semi-supervised learning combines a small amount of labeled data with a large pool of 
unlabeled data during training. This paradigm is particularly relevant in medical 
applications, where expert annotation of data — such as delineating tumor boundaries 
in MRI scans or labeling EEG episodes — is expensive, time-consuming, and requires 
specialized clinical knowledge \cite{chapelle2009semi}.

\paragraph{Reinforcement Learning}
Reinforcement learning (RL) involves an agent that learns optimal decision-making 
policies by interacting with an environment and receiving reward signals. In healthcare, 
RL has been explored for sequential treatment planning, dynamic dosage optimization, 
and personalized therapy recommendation \cite{gottesman2019guidelines}.

% ------------------------------------------------------------
\subsubsection{Common Machine Learning Algorithms}
% ------------------------------------------------------------

Several ML algorithms have been widely adopted in healthcare classification and prediction 
tasks. The most relevant to this work are briefly described below:

\begin{itemize}

    \item \textbf{Logistic Regression}: A probabilistic linear classifier that estimates 
    the probability of class membership using a sigmoid function. Despite its simplicity, 
    it provides interpretable outputs and serves as a useful baseline in clinical settings.

    \item \textbf{Decision Trees}: Hierarchical structures that recursively partition the 
    feature space based on threshold conditions. They are highly interpretable and form 
    the basis for more powerful ensemble methods.

    \item \textbf{Random Forest}: An ensemble of decision trees trained on bootstrapped 
    subsets of the data. By averaging predictions across trees, Random Forest reduces 
    variance and improves generalization, making it effective for high-dimensional 
    biomedical datasets \cite{breiman2001random}.

    \item \textbf{Support Vector Machines (SVM)}: Discriminative classifiers that find 
    the maximum-margin hyperplane separating classes in feature space. Through the kernel 
    trick, SVMs can model non-linear decision boundaries and have been extensively used 
    in EEG-based BCI classification and cancer detection tasks.

    \item \textbf{k-Nearest Neighbors (k-NN)}: A non-parametric classifier that assigns 
    a label based on the majority class among the $k$ nearest samples in feature space. 
    Simple but sensitive to data dimensionality and scaling.

    \item \textbf{XGBoost / Gradient Boosting}: Sequential ensemble methods that build 
    an additive model by fitting each new learner to the residual errors of the previous 
    one. XGBoost in particular has demonstrated strong performance on structured clinical 
    and tabular medical data \cite{chen2016xgboost}.

\end{itemize}

% ------------------------------------------------------------
\subsubsection{Challenges of Machine Learning in Healthcare}
% ------------------------------------------------------------

Despite its promise, deploying ML in healthcare presents several persistent challenges:

\begin{itemize}
    \item \textbf{Data heterogeneity}: Medical data spans multiple modalities — imaging, 
    signals, text, and tabular records — each requiring distinct preprocessing and 
    representation strategies.
    
    \item \textbf{Class imbalance}: In many clinical tasks, such as rare disease detection 
    or adverse event prediction, the minority class is severely underrepresented, leading 
    to biased models.
    
    \item \textbf{Annotation cost}: Labeling medical data requires expert clinical 
    knowledge, making large-scale annotated datasets difficult and costly to produce.
    
    \item \textbf{Interpretability}: Clinical adoption of ML models requires that 
    decisions be explainable to clinicians, creating tension with the opacity of 
    complex models.
    
    \item \textbf{Privacy and data governance}: Patient data is subject to strict 
    regulatory frameworks (GDPR, HIPAA), severely limiting the ability to centralize 
    data across institutions for training — a core motivation for federated 
    learning \cite{miotto2018deep}.
\end{itemize}

% ============================================================
\subsection{Deep Learning}
% ============================================================

\subsubsection{Definition and Principles}

Deep learning (DL) is a subset of machine learning that uses artificial neural networks 
with multiple successive layers to automatically learn hierarchical representations 
from raw data \cite{lecun2015deep}. Rather than relying on manually engineered features, 
DL models discover task-relevant representations directly from input data, which 
has proven transformative for unstructured modalities such as images, audio, and 
biosignals. A single neuron computes:

\[
y = f\!\left(\sum_{i=1}^{n} w_i x_i + b\right)
\]

where $w_i$ are learnable weights, $x_i$ are inputs, $b$ is a bias term, and $f$ 
is a non-linear activation function (e.g., ReLU, sigmoid, or softmax). Stacking 
such neurons in layers enables the network to approximate arbitrarily complex functions.

% ------------------------------------------------------------
\subsubsection{Training Deep Neural Networks}
% ------------------------------------------------------------

DL models are trained by minimizing a loss function $\mathcal{L}$ that measures the 
discrepancy between predictions $\hat{y}$ and ground-truth labels $y$. For 
multi-class classification, the standard choice is the categorical cross-entropy loss:

\[
\mathcal{L} = -\sum_{i=1}^{N} \sum_{c=1}^{C} y_{i,c} \log(\hat{y}_{i,c})
\]

where $N$ is the number of samples, $C$ the number of classes, and $\hat{y}_{i,c}$ 
the predicted probability for class $c$. Weights are updated via backpropagation, 
which computes gradients of $\mathcal{L}$ with respect to each parameter using the 
chain rule. Popular optimizers include:

\begin{itemize}
    \item \textbf{Stochastic Gradient Descent (SGD)}: Updates weights using gradients 
    computed on mini-batches.
    \item \textbf{Adam} \cite{kingma2014adam}: Combines adaptive learning rates with 
    momentum, achieving faster and more stable convergence in most settings.
    \item \textbf{RMSProp}: Adapts learning rates based on a moving average of recent 
    gradient magnitudes, particularly suitable for non-stationary settings.
\end{itemize}

\paragraph{Overfitting and Regularization}
Overfitting occurs when a model memorizes training data but fails to generalize to 
unseen samples. Common mitigation strategies include:
\begin{itemize}
    \item \textbf{L1/L2 regularization}: Penalizes large weight values.
    \item \textbf{Dropout}: Randomly deactivates neurons during training, forcing 
    the network to learn redundant representations.
    \item \textbf{Batch Normalization}: Normalizes layer activations, stabilizing 
    and accelerating training.
    \item \textbf{Early stopping}: Halts training when validation performance 
    stops improving.
\end{itemize}


\subsubsection{Artificial Neural Networks (ANNs) and Multilayer Perceptrons}
The simplest deep learning architecture is the \textbf{multilayer perceptron (MLP)}, 
a fully connected feedforward network in which every neuron in one layer is connected 
to every neuron in the next. An MLP with one hidden layer computes:
\[
\hat{y} = f_2(W_2 \cdot f_1(W_1 \mathbf{x} + b_1) + b_2)
\]
where $f_1$ and $f_2$ are activation functions and $W_1, W_2$ are weight matrices. 
Deeper MLPs can approximate arbitrarily complex functions \cite{hornik1989multilayer}. 
In healthcare, MLPs are commonly applied to structured tabular data such as 
clinical records and laboratory results, and serve as the backbone for more 
specialized architectures such as CNNs and RNNs.

% ------------------------------------------------------------
\subsubsection{Convolutional Neural Networks (CNNs)}
% ------------------------------------------------------------

CNNs are specialized architectures designed for grid-structured data, most notably 
images. Their core components are:

\begin{itemize}
    \item \textbf{Convolutional layers}: Apply learnable filters (kernels) across the 
    input to extract local spatial features such as edges, textures, and shapes.
    \item \textbf{Pooling layers}: Downsample feature maps (e.g., max pooling), 
    reducing spatial dimensions while retaining dominant activations.
    \item \textbf{Fully connected layers}: Aggregate the extracted feature maps to 
    produce final class predictions.
\end{itemize}

A convolutional operation for a 2D input $\mathbf{X}$ with kernel $\mathbf{K}$ is:

\[
(X * K)[i,j] = \sum_{m}\sum_{n} X[i+m,\, j+n] \cdot K[m,n]
\]

CNNs have become the standard approach for medical image analysis, with successful 
applications in tumor detection from MRI and CT scans, retinal disease grading 
from fundus images, and histopathology slide classification \cite{litjens2017survey}. 
Architectures such as ResNet, VGG, and EfficientNet, pre-trained on large natural 
image datasets, are routinely fine-tuned on medical imaging tasks through 
transfer learning.

% ------------------------------------------------------------
\subsubsection{Recurrent Neural Networks (RNNs) and EEGNet}
% ------------------------------------------------------------

RNNs are designed for sequential and temporal data, where the current output depends 
on both the current input and a hidden state encoding prior context:

\[
h_t = f(W_h h_{t-1} + W_x x_t + b)
\]

Vanilla RNNs suffer from vanishing gradients over long sequences. Two gated variants 
address this:

\begin{itemize}
    \item \textbf{Long Short-Term Memory (LSTM)}: Introduces input, forget, and output 
    gates to selectively retain or discard information over time.
    \item \textbf{Gated Recurrent Unit (GRU)}: A simplified gated architecture with 
    fewer parameters, offering comparable performance to LSTM in many settings.
\end{itemize}

RNNs and their variants have been applied to EEG signal analysis, ECG arrhythmia 
detection, and clinical time-series prediction \cite{lipton2015critical}.

For EEG-based classification specifically, \textbf{EEGNet} \cite{lawhern2018eegnet} 
is a compact CNN architecture purpose-built for brain-computer interface (BCI) 
tasks. EEGNet uses depthwise and separable convolutions to learn both temporal 
and spatial EEG filters in a parameter-efficient manner, achieving strong 
generalization across different BCI paradigms with limited training data — a 
key advantage given the small dataset sizes typical in clinical EEG studies. 
Its efficiency makes it particularly well-suited for federated settings where 
local client datasets may be limited.

% ------------------------------------------------------------
\subsubsection{Attention Mechanisms and Transformers}
% ------------------------------------------------------------

Originally introduced for natural language processing, the \textbf{Transformer} 
architecture \cite{vaswani2017attention} has emerged as a powerful paradigm across 
modalities. Its core mechanism, \textbf{scaled dot-product self-attention}, allows 
the model to dynamically weight the relevance of each element in a sequence relative 
to all others:

\[
\text{Attention}(Q, K, V) = \text{softmax}\!\left(\frac{QK^\top}{\sqrt{d_k}}\right)V
\]

where $Q$, $K$, and $V$ are query, key, and value matrices derived from the input, 
and $d_k$ is the key dimension used for scaling. Multi-head attention runs several 
such operations in parallel, capturing relationships at different representation 
subspaces.

In medical imaging, \textbf{Vision Transformers (ViT)} treat image patches as tokens 
and have shown competitive performance with CNNs on classification tasks. In EEG 
analysis, transformer-based models capture long-range temporal dependencies across 
electrode channels, which fixed-kernel CNNs and RNNs may miss. More broadly, 
transformers underpin recent multimodal fusion architectures that jointly process 
imaging, text, and tabular clinical data — a capability increasingly explored in 
federated multimodal healthcare systems.

% ------------------------------------------------------------
\subsubsection{Deep Learning for Multimodal Medical Data}
% ------------------------------------------------------------

A key strength of deep learning is its ability to process diverse medical data 
modalities within unified or jointly trained architectures. The main modality-specific 
approaches are:

\begin{itemize}
    \item \textbf{Medical Imaging} (MRI, CT, X-ray, dermoscopy): CNNs and ViTs are 
    the dominant architectures, with transfer learning from pre-trained models 
    enabling high performance even on small medical datasets.
    \item \textbf{Biosignals} (EEG, ECG, EMG): Temporal models (RNNs, LSTMs) and 
    specialized compact architectures (EEGNet) extract meaningful representations 
    from multi-channel time series.
    \item \textbf{Clinical and tabular data}: Gradient boosting methods (XGBoost) 
    and multilayer perceptrons (MLPs) remain strong performers for structured 
    patient records.
    \item \textbf{Multimodal fusion}: Combining imaging, signals, and tabular data 
    through early, late, or intermediate fusion strategies allows models to 
    leverage complementary information across modalities, improving diagnostic 
    accuracy \cite{miotto2018deep}.
\end{itemize}

This thesis focuses on two of these modalities — EEG signals for epilepsy detection 
and medical images for cancer sub-type classification — both processed through 
federated learning pipelines.

% ============================================================
\subsection{Model Evaluation}
% ============================================================

\subsubsection{Confusion Matrix}

The confusion matrix is the fundamental summary of a classifier's behavior on a 
test set, cross-tabulating predicted labels against true labels. For binary 
classification, it yields four quantities:

\begin{itemize}
    \item \textbf{True Positives (TP)}: Correctly predicted positive cases.
    \item \textbf{True Negatives (TN)}: Correctly predicted negative cases.
    \item \textbf{False Positives (FP)}: Negative cases incorrectly predicted as positive.
    \item \textbf{False Negatives (FN)}: Positive cases incorrectly predicted as negative.
\end{itemize}

% ------------------------------------------------------------
\subsubsection{Classification Metrics}
% ------------------------------------------------------------

From the confusion matrix, the following metrics are derived:

\begin{itemize}
    \item \textbf{Accuracy}: Overall fraction of correct predictions:
    \[
    \text{Accuracy} = \frac{TP + TN}{TP + TN + FP + FN}
    \]
    Accuracy can be misleading under class imbalance, a frequent occurrence in 
    medical datasets.

    \item \textbf{Precision}: Fraction of positive predictions that are truly positive:
    \[
    \text{Precision} = \frac{TP}{TP + FP}
    \]

    \item \textbf{Recall (Sensitivity)}: Fraction of actual positives correctly 
    identified:
    \[
    \text{Recall} = \frac{TP}{TP + FN}
    \]
    In clinical contexts, high recall is often prioritized to minimize missed diagnoses.

    \item \textbf{F1-Score}: Harmonic mean of precision and recall, providing a 
    balanced measure:
    \[
    F_1 = 2 \times \frac{\text{Precision} \times \text{Recall}}{\text{Precision} + \text{Recall}}
    \]

    \item \textbf{Specificity}: Fraction of actual negatives correctly identified:
    \[
    \text{Specificity} = \frac{TN}{TN + FP}
    \]
\end{itemize}

% ------------------------------------------------------------
\subsubsection{ROC Curve and AUC}
% ------------------------------------------------------------

The Receiver Operating Characteristic (ROC) curve plots the True Positive Rate 
(Recall) against the False Positive Rate ($FP / (FP + TN)$) across all 
classification thresholds. The Area Under the Curve (AUC) summarizes this 
in a single scalar: a value of 1.0 indicates perfect discrimination, while 0.5 
corresponds to random chance. AUC is threshold-independent and robust to class 
imbalance, making it a standard evaluation metric in clinical AI research.

% ------------------------------------------------------------
\subsubsection{Convergence Metrics in Federated Settings}
% ------------------------------------------------------------

Beyond standard classification metrics, federated learning introduces additional 
evaluation dimensions. The \textbf{convergence curve} — tracking global model 
accuracy or loss as a function of communication rounds — is essential for 
assessing the stability and efficiency of training across distributed clients. 
Metrics such as \textbf{rounds to convergence} and \textbf{communication cost 
per round} are used alongside task performance to compare aggregation strategies, 
as developed in Chapter~\ref{chap:fl}.

\bigskip

While the AI methods described in this section have demonstrated strong performance 
in centralized settings, their deployment across real-world healthcare institutions 
is fundamentally constrained by data privacy regulations and the practical 
impossibility of pooling sensitive patient data from multiple sites. These 
constraints directly motivate the federated learning paradigm, which is introduced 
and analyzed in depth in the following chapter.


\section{Conclution}
TODO TODO TODO TODO TODO TODO TODO TODO TODO TODO 
However, traditional AI approaches rely on centralized data collection, which is particularly problematic in healthcare due to privacy, legal, and institutional constraints. This has led to the emergence of new paradigms such as federated learning, which enables collaborative model training without sharing raw data. This paradigm will be explored in detail in Chapter 2.






